% !TeX program = pdflatex
% ==========================================================================
%  HPDC-MaOEA -- single-column manuscript draft
%
%  Section 3 (method) is drafted from the two Markdown method drafts
%  ("Hybrid PBI" and "candidate-selection module") and verified line by line
%  against the MATLAB source of
%      REMO_new2_AdaMaO_SDEOnly_UniformMix_Original
%  Sections 1, 2, 4 and 5 are placeholders.
% ==========================================================================
\documentclass[11pt,a4paper]{article}

\usepackage[margin=1in]{geometry}
\usepackage{amsmath,amssymb,amsthm}
\usepackage{bm}
\usepackage{graphicx}
\usepackage{booktabs}
\usepackage{algorithm}
\usepackage{algpseudocode}
\usepackage{enumitem}
\usepackage{xcolor}
\usepackage[hidelinks]{hyperref}

\newtheorem{proposition}{Proposition}
\newtheorem{remark}{Remark}

\newcommand{\vf}{\mathbf{f}}
\newcommand{\vx}{\mathbf{x}}
\newcommand{\vv}{\mathbf{v}}
\newcommand{\vw}{\mathbf{w}}
\newcommand{\vr}{\mathbf{r}}
\newcommand{\vz}{\mathbf{z}}
\newcommand{\keff}{k_{\mathrm{eff}}}
\newcommand{\TODO}[1]{\textcolor{red}{[TODO: #1]}}

\title{Hybrid PBI Quality Stratification and Dual-Mode Candidate Selection
for Expensive Many-Objective Optimization}
\author{Author list to be completed}
\date{\today}

\begin{document}
\maketitle

\begin{abstract}
\TODO{abstract. Reserved memory sentence: across 250 independent runs and 40
predefined Problem--$M$--Stage cells, the two PBI views show low set overlap and
provide bidirectional unique future true positives in 33 cells, and $81.5\%$ of
the union true positives are identified by only one view.}
\end{abstract}

\section{Introduction}
\label{sec:intro}

\TODO{Introduction. Four paragraphs are already drafted in the
problem-mechanism-evidence note (hybrid PBI), Section 5: (i) task and
importance of relation learning under a severely limited evaluation budget;
(ii) the granularity-compression gap of binary PBI partitioning; (iii) the
one-to-one correspondence between the gap and the proposed mechanism;
(iv) strongest evidence and the three contributions.}

\section{Related Work}
\label{sec:related}

\subsection{Surrogate-Assisted Expensive Multi-/Many-Objective Optimization}
\label{sec:related:saea}
\TODO{Kriging/RBF-based SAEAs, their scalability limit in $M$, and why
regression on every objective becomes unattractive when $M$ grows.}

\subsection{Relation Learning for Expensive Optimization}
\label{sec:related:relation}
\TODO{REMO and the pairwise-comparison SAEA line; what supervision signal each
of them uses; position the present work as a labelling-interface contribution
rather than a new scalarizing function.}

\subsection{Quality Assessment and Candidate Selection}
\label{sec:related:selection}
\TODO{Indicator-based assessment (SDE, $L_p$ shape estimation, PIEA), radial
representative selection (RSEA), and infill/candidate selection rules. End with
the gap that motivates Section~\ref{sec:method}.}

\section{The Proposed HPDC-MaOEA}
\label{sec:method}

Relation-based surrogate-assisted evolutionary algorithms replace multi-output
regression by a classifier that predicts the relative quality of a pair of
solutions. Two interfaces then decide what such an algorithm can actually
learn and exploit: the rule that converts the current population into
supervision labels, and the rule that converts the surrogate scores of a large
candidate pool into a small batch of expensive evaluations. HPDC-MaOEA
(hybrid-PBI dual-mode-candidate many-objective evolutionary algorithm)
addresses both interfaces. Section~\ref{sec:method:framework} states the
overall framework, Section~\ref{sec:method:hpc} develops the hybrid PBI quality
stratification that produces the supervision labels,
Section~\ref{sec:method:relation} describes the relation model built on those
labels, Section~\ref{sec:method:candidate} develops the dual-mode candidate
selection, and Section~\ref{sec:method:complexity} analyses the computational
complexity.

\subsection{Overall Framework}
\label{sec:method:framework}

We consider the expensive many-objective problem
\begin{equation}
\min_{\vx\in\Omega}\;
\vf(\vx)=\bigl[f_1(\vx),\ldots,f_M(\vx)\bigr],
\qquad
\Omega\subseteq\mathbb{R}^{D},
\label{eq:problem}
\end{equation}
where every evaluation of $\vf$ is expensive, so the total number of real
function evaluations is limited to a budget $B$. Let $FE$ denote the number of
evaluations consumed so far and
\begin{equation}
\rho=\min\!\left(1,\frac{FE}{B}\right)
\label{eq:rho}
\end{equation}
the evaluation progress.

HPDC-MaOEA starts from a Latin hypercube design of
$N_{\mathrm{init}}$ solutions, evaluates them, and stores them in an archive
$\mathcal{A}_{\mathrm{rc}}$. Each outer iteration then executes five steps.

\begin{enumerate}[leftmargin=*,itemsep=2pt]
\item \textbf{Quality stratification.} The current population $\mathcal{P}$ is
      converted into a binary group label $c_i\in\{0,1\}$ by hybrid PBI
      stratification (Section~\ref{sec:method:hpc}), which also returns the
      anchor set $\mathcal{R}$ used later as an additional mating pool.
\item \textbf{Relation model.} Ordered pairs are generated from $c_i$ and a
      three-class relation network is trained, together with a held-out
      pairwise error rate $p_{\mathrm{err}}$
      (Section~\ref{sec:method:relation}).
\item \textbf{Indicator model.} An SDE fitness value is computed for every
      evaluated solution and an RBF-SVR model is fitted from decision vectors
      to those values (Section~\ref{sec:method:candidate:ind}).
\item \textbf{Candidate selection.} A surrogate-assisted inner search
      accumulates a candidate pool, one of two selection modes is drawn, and a
      batch of $n_t\in[n_{\min},n_{\max}]$ candidates is returned
      (Section~\ref{sec:method:candidate}).
\item \textbf{Evaluation and environmental selection.} The batch is evaluated,
      appended to $\mathcal{A}_{\mathrm{rc}}$, and the next population is
      obtained by radar-grid environmental selection over the whole archive.
\end{enumerate}

Algorithm~\ref{alg:framework} summarises the procedure. The design keeps the
two interfaces separate: hybrid PBI stratification decides \emph{what the
classifier is asked to learn}, whereas dual-mode candidate selection decides
\emph{how its scores are spent}. All expensive evaluations are consumed in
step~5, so both modules operate entirely on already evaluated data and
surrogate predictions.

\begin{algorithm}[t]
\caption{HPDC-MaOEA}
\label{alg:framework}
\begin{algorithmic}[1]
\Require problem, budget $B$, population size $N_p$, candidate-generation limit
         $g_{\max}$, group ratio $r_g$, quantile $q_{\mathrm{keep}}$,
         exploration strength $\lambda_0$, mixing probability
         $p_{\mathrm{mix}}$, batch bounds $n_{\min},n_{\max}$
\Ensure archive $\mathcal{A}_{\mathrm{rc}}$
\State Sample $N_{\mathrm{init}}$ solutions by Latin hypercube design and
       evaluate them
\State $\mathcal{P}\gets$ initial population;
       $\mathcal{A}_{\mathrm{rc}}\gets\mathcal{P}$
\State Initialise the mode stream from the run index and set $L_p\gets1$
\While{$FE<B$}
  \State $\rho\gets FE/B$;
         $\keff\gets\min\bigl(N_p,\max(6,\lceil1.5M\rceil)\bigr)$
  \State $(\bm{c},\mathcal{R})\gets$
         \Call{HybridPBIStratification}{$\mathcal{P},\rho,\keff$}
         \Comment{Section~\ref{sec:method:hpc}}
  \State $(\mathcal{X},\bm{y})\gets$ \Call{RelationPairs}{$\mathcal{P},\bm{c}$}
  \If{$\mathcal{X}=\emptyset$}
    \State $\mathcal{P}\gets$
           \Call{EnvSelect}{$\mathcal{A}_{\mathrm{rc}},N_p$};
           \textbf{continue}
  \EndIf
  \State $(\mathrm{net},p_{\mathrm{err}})\gets$
         \Call{TrainRelationNet}{$\mathcal{X},\bm{y}$}
         \Comment{Section~\ref{sec:method:relation}}
  \State $(\bm{\phi},L_p)\gets$ \Call{SDEFitness}{$\mathcal{P},L_p$};
         $\widehat\phi\gets$ \Call{FitSVR}{$\mathcal{P},\bm{\phi}$}
  \State Draw $\xi\sim U[0,1)$ from the mode stream and resolve the mode
         $m_t$ by \eqref{eq:mode}
  \State $\mathcal{S}\gets$
         \Call{DualModeSelection}{$\mathcal{P},\mathcal{R},\mathrm{net},
         p_{\mathrm{err}},\widehat\phi,m_t,g_{\max}$}
         \Comment{Section~\ref{sec:method:candidate}}
  \If{$\mathcal{S}=\emptyset$}
    \State $\mathcal{S}\gets$ at most $n_{\min}$ offspring generated from
           $\mathcal{P}\cup\mathcal{R}$
  \EndIf
  \State Truncate $\mathcal{S}$ by the remaining budget, evaluate it, and
         append it to $\mathcal{A}_{\mathrm{rc}}$
  \State $\mathcal{P}\gets$ \Call{EnvSelect}{$\mathcal{A}_{\mathrm{rc}},N_p$}
\EndWhile
\end{algorithmic}
\end{algorithm}

\subsection{Hybrid PBI-Based Quality Stratification}
\label{sec:method:hpc}

Relation learning needs a coarse but reliable notion of ``better'' before any
pair can be labelled. A binary PBI partition, as used for relation learning in
REMO~\cite{ref:remo}, provides exactly that: an adaptive boundary divides the
current population into two relatively balanced subpopulations and yields a
clear class structure. The same boundary, however, compresses every PBI
difference inside a subpopulation into a single group-level meaning, so two
solutions of clearly different PBI quality on the same side of the boundary
become indistinguishable in the labels. Hybrid PBI stratification separates the
two roles that the single boundary was carrying. A continuous
direction-based preference preserves the within-group order, an anchor-based
binary preference supplies the between-group boundary, and a stage-aware
combination of the two produces the final positive group used for supervision.

Formally, let the current population be
\begin{equation}
\mathcal{P}=\{\vx_i,\vf_i\}_{i=1}^{N_t},
\qquad
\vf_i\in\mathbb{R}^{M},
\label{eq:pop}
\end{equation}
where $N_t=|\mathcal{P}|$. The module receives the objective matrix, the
progress $\rho$ of \eqref{eq:rho}, and the anchor capacity $\keff$, and returns
a fused sorting score $h_i$, a binary group label $c_i$, and the anchor set
$\mathcal{R}$. Throughout this section $\vz^{*}$ denotes the component-wise
population minimum
\begin{equation}
\vz^{*}=\min_{i=1,\ldots,N_t}\vf_i .
\label{eq:ideal}
\end{equation}
Both branches are computed from the same population and both use PBI geometry;
``hybrid'' therefore refers to the resolution of the two preferences (continuous
versus binary), not to statistical or geometric independence between them. Beyond
resolution the two branches differ in two further implementation respects: the
source of the directions (the continuous branch uses the rays of the nondominated
front itself, \eqref{eq:direction-collapsed}, whereas the binary branch uses the
anchors selected by RSEA, \eqref{eq:anchor-dir}) and the origin of normalisation
(the former is normalised about the origin, the latter about $\vz^{*}$). We do not
claim that these two differences contribute additional informational
complementarity.

\subsubsection{Continuous Direction-Based PBI Preference}
\label{sec:method:hpc:cont}

The continuous branch scores each solution along a direction set
$\mathcal{V}=\{\vv_j\}_{j=1}^{N_v}$, whose size is fixed to the size
$N_{\mathrm{init}}$ of the initial design and therefore equals the population
size in our settings.

\paragraph{Direction set.}
When the number of objectives is small ($M\leq3$) or the population is small
($N_t<50$), $\mathcal{V}$ is a set of uniform reference vectors normalised to
unit length. Otherwise the directions are derived from the population itself.
Let
\begin{equation}
\mathcal{P}^{\mathrm{ND}}
=\bigl\{\vf^{(1)},\ldots,\vf^{(n_{\mathrm{ND}})}\bigr\}
\label{eq:ndset}
\end{equation}
be the first nondominated front of $\mathcal{P}$, and let $\vz_{\mathrm{ND}}$
and $\vz_{\mathrm{ND}}^{\max}$ be its component-wise minimum and maximum, with
range $\bm{\varrho}=\vz_{\mathrm{ND}}^{\max}-\vz_{\mathrm{ND}}$. The front is
first mapped to the unit box,
\begin{equation}
\bar{\vf}^{(j)}
=\bigl(\vf^{(j)}-\vz_{\mathrm{ND}}\bigr)\oslash\bm{\varrho},
\label{eq:ndnorm}
\end{equation}
where $\oslash$ denotes component-wise division. A $K$-means clustering with
$K_v=\min(N_v,n_{\mathrm{ND}})$ clusters is applied to
$\{\bar{\vf}^{(j)}\}$, and every centroid $\bm{\mu}_k$ is mapped back to the
original objective scale and normalised,
\begin{equation}
\vv_k=\frac{\bm{\mu}_k\odot\bm{\varrho}+\vz_{\mathrm{ND}}}
{\bigl\|\bm{\mu}_k\odot\bm{\varrho}+\vz_{\mathrm{ND}}\bigr\|_2},
\qquad k=1,\ldots,K_v,
\label{eq:direction}
\end{equation}
with $\odot$ the component-wise product. If $K_v<N_v$, the centroids are
cyclically replicated until $N_v$ directions are available.

\paragraph{The clustering is an identity map under our settings.}
As stated at the beginning of this subsection, $N_v$ is fixed to the size
$N_{\mathrm{init}}$ of the initial design, whereas $\mathcal{P}^{\mathrm{ND}}$ is
a subset of the current population of size $N_p$, so $n_{\mathrm{ND}}\leq N_p$
always holds. Whenever $N_{\mathrm{init}}\geq N_p$ --- which is the case for all
experiments in this paper, where $D=30$ gives $N_{\mathrm{init}}=N_p=100$ ---
it follows that
\begin{equation}
K_v=\min(N_v,n_{\mathrm{ND}})=n_{\mathrm{ND}},
\label{eq:kvdegen}
\end{equation}
that is, the number of clusters equals the number of samples. The one-point-per-cluster
partition then attains a within-cluster sum of squares of zero, which is the global
optimum, so every centroid coincides with a nondominated solution. Because the affine
normalisation of \eqref{eq:ndnorm} and the back-mapping inside \eqref{eq:direction} are
mutual inverses, \eqref{eq:direction} reduces to
\begin{equation}
\mathcal{V}=\Bigl\{\vf^{(j)}\big/\bigl\|\vf^{(j)}\bigr\|_2
\;\Bigm|\;j=1,\ldots,n_{\mathrm{ND}}\Bigr\},
\label{eq:direction-collapsed}
\end{equation}
the set of \emph{radial unit directions of the current nondominated front itself}:
each direction passes through one observed nondominated solution, the number of
distinct directions equals $n_{\mathrm{ND}}$, and any shortfall with respect to
$N_v$ is filled by cyclic replication. The branch therefore does not summarise
regions of the front by a few centroids; it assigns \emph{one direction per
nondominated solution}. We verified \eqref{eq:kvdegen}--\eqref{eq:direction-collapsed}
term by term on synthetic fronts with $M=5,10,20$ and $N_p=100$: the maximum
deviation between the centroid set and the nondominated point set is $0$, and the
maximum deviation between $\mathcal{V}$ from \eqref{eq:direction} and
\eqref{eq:direction-collapsed} is $2\times10^{-16}$.

The $K$-means call is retained in the implementation so that consumption of the
global random stream matches the completed experiments; as far as the resulting
directions are concerned it can be replaced verbatim by
\eqref{eq:direction-collapsed}. The construction has no shape parameter and
presumes nothing about the convexity or the intrinsic dimension of the front.

Uniform reference vectors are used instead whenever the front is too small to
support clustering, $n_{\mathrm{ND}}<\max(10,N_v/2)$, or whenever some
objective is degenerate over the front,
$\min_m\varrho_m<10^{-12}$; the same fallback is taken if nondominated sorting
or clustering fails. The first condition activates the population-derived field
only when the sample covers the front, and the second avoids a division by
zero in \eqref{eq:ndnorm}. Consequently the \emph{number} of directions follows
the size of the current nondominated set and their \emph{positions} follow the
best solutions obtained so far, so that $\mathcal{V}$ tracks the region the
population currently covers.

\paragraph{Continuous score.}
Each objective vector is associated with the direction of largest cosine
similarity,
\begin{equation}
a(i)=\arg\max_{j}
\frac{\vf_i^{\mathsf T}\vv_j}{\|\vf_i\|_2\,\|\vv_j\|_2},
\label{eq:assoc-v}
\end{equation}
and the parallel and perpendicular distances with respect to that direction are
\begin{align}
d_{1,i}&=\frac{(\vf_i-\vz^{*})^{\mathsf T}\vv_{a(i)}}{\|\vv_{a(i)}\|_2},
\label{eq:d1}\\[2pt]
d_{2,i}&=\left\|(\vf_i-\vz^{*})
-d_{1,i}\frac{\vv_{a(i)}}{\|\vv_{a(i)}\|_2}\right\|_2 .
\label{eq:d2}
\end{align}
The continuous PBI value follows the conventional penalty-based boundary
intersection form~\cite{ref:moead}, and its monotone score is obtained by a
decreasing transform,
\begin{equation}
g_i^{\mathrm{con}}=d_{1,i}+\theta\,d_{2,i},
\qquad
s_i=\frac{1}{1+g_i^{\mathrm{con}}},
\label{eq:score-cont}
\end{equation}
with penalty parameter $\theta=5$. A larger $s_i$ indicates a smaller PBI value
along the associated direction, and unlike a hard class label $s_i$ retains a
fine-grained ordering of the current population.

Equations \eqref{eq:assoc-v}--\eqref{eq:d2} also fix the coordinate convention
of the implementation: the association is computed from the original objective
vectors, whereas the two distances are computed after translation by the ideal
point $\vz^{*}$. All experiments in this paper use this single convention.

\paragraph{Closed form of the continuous score under the collapsed direction set.}
Combining \eqref{eq:direction-collapsed} with \eqref{eq:assoc-v} has one
consequence that must be made explicit. Since $\mathcal{V}$ consists of the rays
of the nondominated solutions themselves, any nondominated $\vf_i$ attains the
maximum value $1$ of \eqref{eq:assoc-v} at the direction generated by itself, and
is therefore associated with $\vv_{a(i)}=\vf_i/\|\vf_i\|_2$ (if collinear
solutions exist, the associated direction is collinear with $\vf_i$ and the
conclusion is unchanged). Writing $\psi_i$ for the angle between $\vf_i$ and
$\vz^{*}$, \eqref{eq:d1}--\eqref{eq:score-cont} give
\begin{align}
d_{1,i}&=\|\vf_i\|_2-\|\vz^{*}\|_2\cos\psi_i,
\qquad
d_{2,i}=\|\vz^{*}\|_2\sin\psi_i,
\label{eq:d-collapsed}\\[2pt]
s_i&=\Bigl(1+\|\vf_i\|_2-\|\vz^{*}\|_2\cos\psi_i
+\theta\,\|\vz^{*}\|_2\sin\psi_i\Bigr)^{-1}.
\label{eq:score-collapsed}
\end{align}
The only solution-dependent leading term in \eqref{eq:score-collapsed} is
$\|\vf_i\|_2$; the remaining two terms enter only through the angle of $\vf_i$
relative to the fixed vector $\vz^{*}$. On nondominated solutions the continuous
branch therefore measures \emph{mainly convergence along the ray rather than the
position of a solution on the front}: it supplies none of the diversity pressure
that decomposition-based algorithms obtain from directions competing against one
another, because no two solutions ever contend for the same direction. For
dominated solutions the associated direction is the nearest nondominated ray and
\eqref{eq:d-collapsed}--\eqref{eq:score-collapsed} no longer hold exactly. Since
the large majority of the population is nondominated once $M\geq10$,
\eqref{eq:score-collapsed} covers most individuals of a given generation.
Section~\ref{sec:exp:hpc} reports the measured rank correlation between $s_i$ and
$-\|\vf_i-\vz^{*}\|_2$ to quantify the extent of this collapse.

\TODO{Export the \texttt{DirectionSource}, \texttt{ClusterCount},
$n_{\mathrm{ND}}$ and $\mathrm{corr}(s_i,-\|\vf_i-\vz^{*}\|_2)$ audit fields once
under the formal protocol ($M\in\{3,5,10,20\}$, $D=30$, $N_p=100$, $B=500$) and
replace the synthetic-front verification numbers of this subsection by real run
logs.}

\subsubsection{Anchor-Based Coarse Preference}
\label{sec:method:hpc:anchor}

The continuous score orders the population along the population-derived
directions but does not provide a class boundary. The second branch supplies
that boundary from a small set of anchors
\begin{equation}
\mathcal{R}=\{\vr_j\}_{j=1}^{\keff}\subset\mathcal{P},
\label{eq:anchors}
\end{equation}
selected by the radial-projection-based representative selection of
RSEA~\cite{ref:rsea}, which is also the selection rule used by
REMO~\cite{ref:remo}. The anchors are evaluated solutions of the current
population and are used as dynamic reference representatives; they are not
treated as samples of the true Pareto front.

\paragraph{Anchor capacity and grid resolution.}
The effective number of anchors grows with the number of objectives,
\begin{equation}
\keff=\min\Bigl(N_p,\;\max\bigl(6,\lceil1.5M\rceil\bigr)\Bigr),
\label{eq:keff}
\end{equation}
so that $M=3,5,10,20$ give $\keff=6,8,15,30$ when $N_p$ is not active as an
upper bound, with six anchors as a lower bound for few objectives. In the
present implementation $\keff$ has a second effect. Representative selection
maps the $M$-dimensional normalised objective vectors onto two-dimensional
radar coordinates and partitions that plane into an
$n_{\mathrm{div}}\times n_{\mathrm{div}}$ grid with
\begin{equation}
n_{\mathrm{div}}=\bigl\lceil\sqrt{\keff}\,\bigr\rceil ,
\label{eq:ndiv}
\end{equation}
then distributes the representatives over the occupied cells, preferring sparse
cells and, inside a cell, trading a small normalised objective sum against a
large radar-projection distance to the already selected solutions. Hence
$\keff$ is simultaneously the capacity of the anchor set and the resolution of
the grid over which anchors are spread: $\keff=6$ yields a $3\times3$ grid and
$\keff=30$ a $6\times6$ grid. This coupling is the direct reason for scaling
\eqref{eq:keff} with $M$; a fixed anchor budget would also fix the grid
resolution, and the anchors could not cover additional radial regions as $M$
grows.

\paragraph{Binary preference.}
Every solution is associated with the anchor of largest cosine similarity in the
original objective space, giving the index $b(i)$ and the unit direction
\begin{equation}
\vw_{b(i)}=\frac{\vr_{b(i)}-\vz^{*}}{\|\vr_{b(i)}-\vz^{*}\|_2}.
\label{eq:anchor-dir}
\end{equation}
The corresponding distances are
\begin{align}
\hat d_{1,i}&=(\vf_i-\vz^{*})^{\mathsf T}\vw_{b(i)},
\label{eq:d1hat}\\[2pt]
\hat d_{2,i}&=\bigl\|(\vf_i-\vz^{*})-\hat d_{1,i}\vw_{b(i)}\bigr\|_2 ,
\label{eq:d2hat}
\end{align}
and, for a balancing parameter $\delta$, the anchor-normalised PBI value and the
binary preference are
\begin{align}
g_i^{\mathrm{bin}}(\delta)
&=\frac{\hat d_{1,i}+\delta\,\hat d_{2,i}}{\|\vr_{b(i)}-\vz^{*}\|_2},
\label{eq:gbin}\\[2pt]
\ell_i(\delta)
&=\mathbb{I}\bigl[g_i^{\mathrm{bin}}(\delta)\leq1\bigr].
\label{eq:label-bin}
\end{align}
Normalising by $\|\vr_{b(i)}-\vz^{*}\|_2$ makes the threshold in
\eqref{eq:label-bin} the anchor itself, so $\ell_i$ records on which side of the
surface passing through the associated anchor a solution lies.

The parameter $\delta$ is a signed class-balance variable rather than a new
scalarizing parameter. REMO adjusts its partition parameter towards an
approximately balanced split; the implementation used here searches
$\delta\in[-20,20]$ by bisection and stops as soon as the positive rate
\begin{equation}
r(\delta)=\frac{1}{N_t}\sum_{i=1}^{N_t}\ell_i(\delta)
\label{eq:posrate}
\end{equation}
enters the tolerant interval $[0.3,0.7]$, or when the bracket width falls below
$10^{-1}$. The bounded search therefore seeks a tolerant class ratio and
terminates in a bounded number of steps; because both the interval and the
tolerance are finite, it does not guarantee a ratio inside $[0.3,0.7]$ for every
population state. A negative $\delta$ is admissible and rewards, rather than
penalises, a large perpendicular distance, which is what allows the same rule to
raise the positive rate when the associated anchors are close to the population.

\subsubsection{Stage-Aware Hybrid Stratification}
\label{sec:method:hpc:hybrid}

The two preferences are combined with a weight that follows the consumption of
the expensive budget. With $\rho$ from \eqref{eq:rho}, the continuous branch
receives the weight
\begin{equation}
\alpha=1-\rho,
\label{eq:alpha}
\end{equation}
and the fused sorting score is
\begin{equation}
h_i=\alpha\,s_i+(1-\alpha)\,\ell_i .
\label{eq:fuse}
\end{equation}
Solutions are sorted in descending order of $h_i$ and the positive group is the
top fraction $r_g$,
\begin{equation}
c_i=
\begin{cases}
1, & \operatorname{rank}_{\downarrow}(h_i)\leq\lceil r_gN_t\rceil,\\[2pt]
0, & \text{otherwise},
\end{cases}
\label{eq:group}
\end{equation}
with $r_g=0.25$. All solutions with $c_i=0$, both mid-ranked and low-ranked
ones, form a single non-positive group.

Equation \eqref{eq:fuse} is an annealing schedule over the sorting key rather
than a plain weighted average of two comparable quantities. Since
$s_i\in(0,1]$ is continuous while $\ell_i\in\{0,1\}$, the two branches act on
different resolvable scales. For large $\alpha$ the differences in $h_i$ are
governed by $s_i$, and the continuous score is the sorting key. Once $\alpha$
is small enough that the gap $1-\alpha$ contributed by the binary preference
exceeds the largest attainable difference in the continuous score, the ordering
changes qualitatively: the binary class becomes the key and $s_i$ degenerates
into a tie-breaker inside each class. Proposition~\ref{prop:twolevel} states a
sufficient condition for that transition, which depends only on the ranges of
$\alpha$ and $s_i$ and requires no assumption on the geometry of the front. The
role of \eqref{eq:alpha} is therefore to move the \emph{granularity} of the
supervision signal monotonically from fine to coarse: the early budget keeps
within-group resolvability, and the late budget enforces a stable class
structure. Section~\ref{sec:exp:hpc} tests this schedule against static and
reversed alternatives.

\begin{proposition}[Two-level ordering induced by the binary preference]
\label{prop:twolevel}
Assume the directions are nonnegative, $\vv_j\succeq\bm{0}$, and the ideal point
is the component-wise population minimum \eqref{eq:ideal}, so that
$\vf_i-\vz^{*}\succeq\bm{0}$. Then $d_{1,i}\geq0$ and $d_{2,i}\geq0$, hence
$g_i^{\mathrm{con}}\geq0$ and $s_i\in(0,1]$ by \eqref{eq:score-cont}. Under
this premise, for any pair $i,j$ with $\ell_i=1$ and $\ell_j=0$,
\begin{equation*}
\alpha\leq\tfrac{1}{2}
\quad\Longrightarrow\quad
h_i>h_j .
\end{equation*}
\end{proposition}

\begin{proof}
By \eqref{eq:fuse}, $h_i=(1-\alpha)+\alpha s_i>1-\alpha$ because $s_i>0$, while
$h_j=\alpha s_j\leq\alpha$ because $s_j\leq1$. For $\alpha\leq1/2$ we have
$1-\alpha\geq\alpha$, and therefore $h_i>h_j$.
\end{proof}

Proposition~\ref{prop:twolevel} shows that as soon as the binary branch holds at
least half of the weight, \eqref{eq:fuse} induces an explicit two-level order:
every binary-positive solution precedes every binary-negative one, and the
continuous score orders solutions within each class. The condition can be
relaxed. Let
\begin{equation}
\Delta_{01}=\max_{j:\ell_j=0}s_j-\min_{i:\ell_i=1}s_i
\label{eq:delta01}
\end{equation}
be the cross-class span of the continuous score. Strict cross-class dominance is
retained whenever
\begin{equation}
\alpha<\frac{1}{1+\Delta_{01}} .
\label{eq:relaxed}
\end{equation}
Because $s_i\in(0,1]$ implies $\Delta_{01}\leq1$ and hence
$1/(1+\Delta_{01})\geq1/2$, the bound $\alpha\leq1/2$ is only the worst case:
whenever the continuous scores of one generation span a narrow range, the binary
class becomes the sorting key earlier than $\rho=1/2$, so the schedule in
\eqref{eq:alpha} extends the binary-dominated regime beyond half of the budget
in practice.

\begin{remark}[Realised weight range]
\label{rem:alphafirst}
Since the initial design is already evaluated when the module is first invoked,
the first realised weight is not $\alpha=1$ but
\begin{equation}
\alpha_{\mathrm{first}}=1-\frac{N_{\mathrm{init}}}{B} .
\label{eq:alphafirst}
\end{equation}
Under the protocol of this paper ($N_{\mathrm{init}}=100$, $B=500$),
$\alpha_{\mathrm{first}}=0.8$, so the sufficient condition $\alpha\leq1/2$ is met
from $FE=250$ onwards. Both the nominal schedule and the weight range actually
visited by the algorithm are reported in Section~\ref{sec:exp:hpc}.

The relaxed condition \eqref{eq:relaxed} makes the actual transition earlier.
Since $\Delta_{01}$ cannot exceed the full span of $s_i$, which is $0.035$--$0.172$
(for $M=20$ down to $M=5$) on the synthetic fronts we verified, substituting into
\eqref{eq:relaxed} places the moment at which the binary branch becomes the
sorting key at $\rho\approx0.03$--$0.15$, that is, shortly after the initial design
has been evaluated. \emph{The schedule of \eqref{eq:alpha} is therefore not two
signals acting jointly over a long horizon, but a continuous branch that serves as
the sorting key only during a small initial fraction of the budget, after which the
binary class takes over and the continuous score degenerates to a within-class
tie-breaker.} This quantitative statement must be read together with the
qualitative description in Section~\ref{sec:method:hpc:hybrid};
Section~\ref{sec:exp:hpc} reports the distribution of $\Delta_{01}$ and the measured
transition point on real runs.
\end{remark}

Algorithm~\ref{alg:hpc} summarises the stratification module.

\begin{algorithm}[t]
\caption{Hybrid PBI quality stratification}
\label{alg:hpc}
\begin{algorithmic}[1]
\Require population $\mathcal{P}$, progress $\rho$, anchor capacity $\keff$,
         penalty $\theta$, group ratio $r_g$
\Ensure group labels $\bm{c}$, fused scores $\bm{h}$, anchors $\mathcal{R}$
\State Build $\mathcal{V}$: uniform vectors if $M\leq3$, $N_t<50$ or the
       fallback conditions of Section~\ref{sec:method:hpc:cont} hold; otherwise
       the radial unit directions of the nondominated front,
       \eqref{eq:direction-collapsed}, cyclically replicated to $N_v$ entries
\State $\vz^{*}\gets$ component-wise minimum of the objectives
\State Associate every solution by \eqref{eq:assoc-v} and compute $s_i$ by
       \eqref{eq:d1}--\eqref{eq:score-cont}
\State $\mathcal{R}\gets$ radial-grid selection of $\keff$ anchors from
       $\mathcal{P}$
\State Bisect $\delta\in[-20,20]$ until $r(\delta)\in[0.3,0.7]$ or the bracket
       is smaller than $10^{-1}$; set $\ell_i$ by \eqref{eq:label-bin}
\State $\alpha\gets1-\rho$;
       $h_i\gets\alpha s_i+(1-\alpha)\ell_i$
\State Mark the $\lceil r_gN_t\rceil$ solutions of largest $h_i$ as the positive
       group and all others as the non-positive group
\State \Return $\bm{c},\bm{h},\mathcal{R}$
\end{algorithmic}
\end{algorithm}

\subsection{Relation Model Construction}
\label{sec:method:relation}

The group labels of \eqref{eq:group} define the supervision target of the
relation model. Following the relation-modelling protocol of
REMO~\cite{ref:remo}, let
\begin{equation}
\mathcal{C}_1=\{\vx_i\mid c_i=1\},
\qquad
\mathcal{C}_2=\{\vx_i\mid c_i\neq1\}
\label{eq:groups}
\end{equation}
be the positive and non-positive groups. Every ordered pair of distinct
solutions is turned into one training sample whose input is the concatenation
$[\vx_i,\vx_j]\in\mathbb{R}^{2D}$ and whose label
\begin{equation}
y_{ij}=
\begin{cases}
+1, & \vx_i\in\mathcal{C}_1,\ \vx_j\in\mathcal{C}_2,\\
-1, & \vx_i\in\mathcal{C}_2,\ \vx_j\in\mathcal{C}_1,\\
0,  & \text{$\vx_i$ and $\vx_j$ in the same group},
\end{cases}
\label{eq:pairlabel}
\end{equation}
records the ordered group relation rather than a Pareto comparison. Self-pairs
are removed. Since $r_g=0.25$ makes the non-positive group three times larger
than the positive one, the two same-group families are subsampled towards the
cross-group ones: with $n_{\times}=\lceil|\mathcal{C}_1\times\mathcal{C}_2|/2
\rceil$ as the target, positive--positive and non-positive--non-positive pairs
are each randomly reduced to $n_{\times}$ samples when both exceed it, and when
one family falls short the other absorbs the deficit so that the total number of
same-group samples remains $2n_{\times}$. The resulting label distribution is
balanced by construction rather than by a reweighting of the loss.

The pair set is split by stratified sampling, $75\%$ for training and $25\%$ as a
held-out set $\mathcal{T}$, so that all three labels of \eqref{eq:pairlabel} are
represented in both parts. Note that the split is performed over pairs and not
over base solutions, so a solution may appear on both sides of the split. Inputs
are min--max normalised, and a feed-forward pattern network whose three hidden
layers have $3D$, $2D$ and $D$ units, i.e.\ $1.5$, $1$ and $0.5$ times the input
dimension $2D$, maps the normalised pair to a softmax output
\begin{equation}
\bm{\pi}(\vx_a,\vx_b)
=\bigl[\pi_{+1},\pi_0,\pi_{-1}\bigr],
\qquad
\pi_{+1}+\pi_0+\pi_{-1}=1 ,
\label{eq:netout}
\end{equation}
whose components are the predicted probabilities that the first solution belongs
to the higher group, that both belong to the same group, and that the second
belongs to the higher group. The held-out misclassification rate
\begin{equation}
p_{\mathrm{err}}
=\frac{1}{|\mathcal{T}|}
\sum_{(a,b)\in\mathcal{T}}
\mathbb{I}\bigl[\hat y_{ab}\neq y_{ab}\bigr]
\label{eq:perr}
\end{equation}
is retained and reused in Section~\ref{sec:method:candidate:exp} to modulate the
exploration term; it is set to one whenever the held-out set is empty or the
estimate is not finite. If the pair set is empty in some iteration, the model is
not trained, the iteration performs environmental selection only, and no
expensive evaluation is spent.

\subsection{Dual-Mode Candidate Selection}
\label{sec:method:candidate}

The relation network ranks a large number of unevaluated solutions, but the
distribution of its output score changes with the problem, the search stage and
the current training set. Applying a fixed absolute cutoff transfers that scale
variation directly to the number of expensive evaluations spent per iteration.
Two facts make the effect concrete for the score used here. First, the score is
analytically bounded, so a conventional cutoff of $3.9$ sits at $97.5\%$ of its
theoretical maximum (Section~\ref{sec:method:candidate:score}). Second, in a
400-run controlled audit over four variants, five problems and 20 independent
runs, the diagnostic arm that recorded but did not apply that cutoff showed a
$5.1$-fold cross-problem variation in the passing proportion; had the cutoff
driven the batch size, the same 200 additional evaluations would have produced
between $14.8$ and $46.2$ surrogate-assisted iterations, a $3.1$-fold
difference. HPDC-MaOEA therefore separates ranking from expenditure: quantile
rules convert scores of any scale into within-pool ranks, and an explicit
interval $[n_{\min},n_{\max}]$ bounds the batch
(Section~\ref{sec:method:candidate:mix}).

\subsubsection{Pairwise Relation Score and Candidate Pool}
\label{sec:method:candidate:score}

\paragraph{Score.}
For an unevaluated candidate $\vx$, four families of ordered comparisons against
the two groups of \eqref{eq:groups} are constructed,
\begin{equation}
(\mathcal{C}_1,\vx),\quad
(\vx,\mathcal{C}_1),\quad
(\mathcal{C}_2,\vx),\quad
(\vx,\mathcal{C}_2),
\label{eq:fourfamilies}
\end{equation}
which amounts to $2(|\mathcal{C}_1|+|\mathcal{C}_2|)=2N_t$ network inputs per
candidate. Let $\mathbf{a}(\vx)$, $\mathbf{b}(\vx)$, $\mathbf{c}(\vx)$ and
$\mathbf{d}(\vx)$ be the family-wise mean probability vectors of
\eqref{eq:netout}. Aggregating the evidence that $\vx$ is not inferior to the
positive group and is superior to the non-positive group, and subtracting the
opposite evidence, gives the pairwise relation score
\begin{equation}
r(\vx)=2\bigl[
c_{-1}(\vx)+d_{+1}(\vx)-a_{+1}(\vx)-b_{-1}(\vx)
\bigr].
\label{eq:relscore}
\end{equation}
Because each probability vector of \eqref{eq:netout} sums to one, every bracket
term lies in $[0,1]$ and
\begin{equation}
-4\leq r(\vx)\leq 4 ,
\label{eq:relbound}
\end{equation}
which is the bound quoted above. The score is a heuristic net evidence relative
to the two coarse quality groups, not a Pareto win rate.

\paragraph{Candidate pool.}
Candidates are produced by a relation-guided inner evolutionary search. Genetic
variation is first applied to the union of the current population and the
anchors, yielding $\mathcal{Q}^{(0)}$. At inner iteration $\ell$ the network
scores $\mathcal{Q}^{(\ell)}$ by \eqref{eq:relscore}, the best
\begin{equation}
n_{\mathrm{parent}}=\min\bigl(|\mathcal{R}|,|\mathcal{Q}^{(\ell)}|\bigr)
\label{eq:nparent}
\end{equation}
candidates are kept, and variation on those parents together with the anchors
produces $\mathcal{Q}^{(\ell+1)}$. The loop stops when the cumulative number of
generated candidates reaches the limit $g_{\max}$, and all candidates generated
along the way are accumulated and deduplicated,
\begin{equation}
\mathcal{A}=\operatorname{unique}
\Bigl(\textstyle\bigcup_{\ell=0}^{L}\mathcal{Q}^{(\ell)}\Bigr).
\label{eq:pool}
\end{equation}
Both selection modes rank the pool $\mathcal{A}$ produced by \eqref{eq:pool} with
their own criterion. The pool is not, however, mode-independent: the mode is drawn
before the inner search starts, and the parent retention of \eqref{eq:nparent} uses
the family aggregation of the active mode --- sharpness-weighted in the exploratory
mode, plain in the indicator mode (Section~\ref{sec:method:candidate:exp}). The
inner trajectory therefore differs between modes, and so does the realised
$\mathcal{A}$. What the two modes share is the pool \emph{construction rule}
\eqref{eq:pool} and the batch bound, not one identical realised candidate set.
Section~\ref{sec:exp:candidate} isolates the ranking criterion from this
confound by re-running both criteria on a shared frozen pool.

\subsubsection{Relation-Guided Exploratory Mode}
\label{sec:method:candidate:exp}

The exploratory mode combines the relation score with the ambiguity of the
classifier output and with distance in decision space.

Let $\mathcal{O}(\vx)$ collect all ordered comparisons of \eqref{eq:fourfamilies}
for candidate $\vx$. The mean classification ambiguity is
\begin{equation}
u(\vx)=1-\frac{1}{|\mathcal{O}(\vx)|}
\sum_{(\vx_a,\vx_b)\in\mathcal{O}(\vx)}
\max_{y\in\{-1,0,+1\}}\pi_y(\vx_a,\vx_b),
\label{eq:ambiguity}
\end{equation}
so a large $u(\vx)$ means the network produces flatter class distributions around
that candidate. Equation~\eqref{eq:ambiguity} describes how concentrated the
softmax output is; it is not a calibrated variance or an epistemic uncertainty.
In this mode the same sharpness $\max_y\pi_y$ is also used as the weight of the
family-wise averages entering \eqref{eq:relscore}, so that sharper pairs
contribute more to the aggregated score. The weighting changes only the
aggregation inside a family and leaves the semantics of the three relation
classes unchanged.

With $\widetilde r$ and $\widetilde u$ denoting min--max normalisation over the
pool, the exploratory score is
\begin{equation}
a^{\mathrm{E}}(\vx)=\widetilde r(\vx)+\lambda_t\,\widetilde u(\vx),
\label{eq:expscore}
\end{equation}
where
\begin{equation}
\lambda_t=\lambda_0\left(1-\frac{FE}{B}\right)
\max\!\left(0,\,1-\frac{p_{\mathrm{err}}}{0.45}\right)
\label{eq:lambda}
\end{equation}
and $\lambda_0=0.35$. The two factors of \eqref{eq:lambda} tie the strength of
ambiguity sampling to the two quantities that make it meaningful: the remaining
budget, and the held-out error \eqref{eq:perr} of the model whose ambiguity is
being sampled. The term vanishes once $p_{\mathrm{err}}\geq0.45$, so a poorly
performing classifier cannot steer the batch through its own flat outputs.

Candidates above the $q_{\mathrm{keep}}$ quantile of the exploratory score form
the preliminary set
\begin{equation}
\mathcal{H}^{\mathrm{E}}=\Bigl\{\vx\in\mathcal{A}\;\Bigm|\;
a^{\mathrm{E}}(\vx)\geq
Q_{q_{\mathrm{keep}}}\bigl(a^{\mathrm{E}}(\mathcal{A})\bigr)\Bigr\},
\label{eq:prefilter-exp}
\end{equation}
with $q_{\mathrm{keep}}=0.80$, which retains roughly the best $20\%$ of the pool;
if fewer than $n_{\min}$ candidates qualify, the set is completed with the
highest-scoring candidates. The batch is then assembled greedily. The first
member maximises $a^{\mathrm{E}}$, and each subsequent member maximises
\begin{equation}
A(\vx\mid\mathcal{S})
=0.75\,\widehat a^{\mathrm{E}}(\vx)+0.25\,\widehat d(\vx,\mathcal{S}),
\qquad
d(\vx,\mathcal{S})=\min_{\vz\in\mathcal{S}}\|\vx-\vz\|_2 ,
\label{eq:greedy}
\end{equation}
where $\widehat a^{\mathrm{E}}$ and $\widehat d$ are normalised over the
remaining candidates at each step. The $0.75$--$0.25$ split keeps relation
quality as the primary criterion while preventing one batch from collapsing into
a single region of decision space. The rule promotes spread in decision space and
does not by itself guarantee spread in objective space.

\subsubsection{Indicator-Guided Mode}
\label{sec:method:candidate:ind}

The indicator-guided mode reranks candidates by a second criterion whose training
target does not come from the relation network. Its fitness is the shift-based
density estimation (SDE) fitness~\cite{ref:sde}, computed with the generalised
$L_p$ front-shape parameter of PIEA~\cite{ref:piea}.

\paragraph{Front-shape estimation.}
The first nondominated front of the evaluated population is normalised to the
unit box, and for each of $17$ candidate exponents
$p\in[0.27,6.5]$ the generalised norms
$G_p(\bar{\vf}_i)=(\sum_m\bar f_{i,m}^{\,p})^{1/p}$ are computed, trimmed by a
boxplot rule with factor $1.5$, and scaled by their maximum. The exponent
attaining the smallest standard deviation is retained as $L_p$, since a matching
exponent makes the front an approximate level surface of that norm. If the front
holds fewer than $20$ solutions, $L_p=1$ is used.

\paragraph{SDE fitness.}
For normalised objective vectors the shifted nearest-neighbour distance is
\begin{equation}
\delta_i=\min_{j\neq i}
\bigl\|\bar{\vf}_i-\max(\bar{\vf}_i,\bar{\vf}_j)\bigr\|_2 ,
\label{eq:sde}
\end{equation}
where the maximum is component-wise, so the distance vanishes only when $\vx_j$
is at least as good as $\vx_i$ in every objective. The values are rescaled to
$[0,3]$; solutions whose rescaled value falls below $10^{-4}$, i.e.\ those the
density estimate cannot discriminate, take instead the negatively rescaled
Minkowski-$L_p$ distance to the current ideal point, so that convergence
information replaces the degenerate density value. A $\tanh$ transform maps the
result to $[-1,1]$, giving the fitness $\phi_i$ of every evaluated solution.

\paragraph{Reranking.}
Since $\phi_i$ is defined only for evaluated solutions, it is extended to the
pool by an RBF-SVR model fitted on the evaluated decision vectors,
\begin{equation}
\widehat\phi(\vx)=\operatorname{SVR}
\bigl(\vx;\mathcal{P},\bm{\phi}\bigr),
\label{eq:svr}
\end{equation}
with automatic kernel scale and standardised inputs. Selection then proceeds in
two stages: the relation score first retains the top $30\%$ of the pool, and at
least $20$ candidates, after which \eqref{eq:svr} reranks the retained set and
candidates not below the $70\%$ quantile of $\widehat\phi$ form
$\mathcal{H}^{\mathrm{I}}$, again completed with the best candidates if fewer
than $n_{\min}$ qualify. The batch is the highest-ranked $n_t$ candidates of that
set under $\widehat\phi$. The ordering of the two stages assigns each model the
task it is trained for: the relation model excludes evidently low-priority
regions, and the SDE surrogate expresses front-shape and local-density
preferences inside the reduced set. If the SVR is unavailable or its predictions
are not finite, the mode falls back to the relation score.

\subsubsection{Mixed-Mode Allocation}
\label{sec:method:candidate:mix}

Let $m_t\in\{\mathrm{E},\mathrm{I}\}$ be the mode used at outer iteration $t$. A
random stream seeded from the run index, and independent of the global stream
that drives variation and subsampling, produces $\xi_t\sim U[0,1)$, and
\begin{equation}
m_t=
\begin{cases}
\mathrm{I}, & \text{if the indicator model is available and }
              \xi_t<p_{\mathrm{mix}},\\[2pt]
\mathrm{E}, & \text{otherwise},
\end{cases}
\label{eq:mode}
\end{equation}
with $p_{\mathrm{mix}}=0.5$. The two ranking criteria thus receive equal prior
exposure without introducing a further calibration model, and the dedicated
stream makes the mode sequence of a run reproducible independently of how many
random numbers the rest of the iteration consumes. The modes are not fused
numerically: each iteration uses exactly one ranking criterion, under the same
pool-construction rule \eqref{eq:pool} and the same batch bound. If no indicator
model could be fitted, the exploratory mode is used.

Whichever mode is drawn, the number of expensive evaluations is
\begin{equation}
n_t=\min\Bigl[n_{\max},\,
\max\bigl(n_{\min},|\mathcal{H}^{m_t}|\bigr)\Bigr],
\label{eq:batch}
\end{equation}
truncated further by the remaining budget $B-FE$, with $n_{\min}=4$ and
$n_{\max}=6$. Together with the quantile rules
\eqref{eq:prefilter-exp} and the $70\%$ rule of
Section~\ref{sec:method:candidate:ind}, \eqref{eq:batch} completes the separation
announced at the beginning of this section: surrogate scores decide the priority
of candidates inside a batch, and no longer decide how much of the expensive
budget one outer iteration consumes. If the selected set is empty, one round of
genetic variation on the population and the anchors supplies at most $n_{\min}$
candidates, so the search loop never stalls.

Algorithm~\ref{alg:candidate} summarises the module.

\begin{algorithm}[t]
\caption{Dual-mode candidate selection}
\label{alg:candidate}
\begin{algorithmic}[1]
\Require population $\mathcal{P}$, anchors $\mathcal{R}$, relation network,
         held-out error $p_{\mathrm{err}}$, progress $FE/B$, limit $g_{\max}$,
         mode draw $\xi$, parameters
         $q_{\mathrm{keep}},\lambda_0,p_{\mathrm{mix}},n_{\min},n_{\max}$
\Ensure candidate batch $\mathcal{S}$
\State $\mathcal{Q}^{(0)}\gets$ variation on $\mathcal{P}\cup\mathcal{R}$;
       $\mathcal{A}\gets\mathcal{Q}^{(0)}$; $\ell\gets0$
\While{the number of generated candidates is below $g_{\max}$ and
       $\mathcal{Q}^{(\ell)}\neq\emptyset$}
  \State Score $\mathcal{Q}^{(\ell)}$ by \eqref{eq:relscore} and keep the best
         $n_{\mathrm{parent}}$ of \eqref{eq:nparent} as parents
  \State $\mathcal{Q}^{(\ell+1)}\gets$ variation on the parents and
         $\mathcal{R}$;
         $\mathcal{A}\gets\operatorname{unique}
         (\mathcal{A}\cup\mathcal{Q}^{(\ell+1)})$; $\ell\gets\ell+1$
\EndWhile
\State Resolve the mode $m_t$ by \eqref{eq:mode}
\If{$m_t=\mathrm{I}$}
  \State Retain the top $30\%$ of $\mathcal{A}$ by \eqref{eq:relscore}
         (at least $20$ candidates)
  \State Rerank them by \eqref{eq:svr} and keep those above its $70\%$ quantile
         as $\mathcal{H}^{\mathrm{I}}$
  \State $\mathcal{S}\gets$ the $n_t$ best candidates of
         $\mathcal{H}^{\mathrm{I}}$ under $\widehat\phi$
\Else
  \State Compute $a^{\mathrm{E}}$ by \eqref{eq:ambiguity}--\eqref{eq:lambda} and
         form $\mathcal{H}^{\mathrm{E}}$ by \eqref{eq:prefilter-exp}
  \State $\mathcal{S}\gets$ greedy batch of $n_t$ candidates by
         \eqref{eq:greedy}
\EndIf
\State Bound $n_t$ by \eqref{eq:batch} and by the remaining budget
\State \Return $\mathcal{S}$
\end{algorithmic}
\end{algorithm}

\subsection{Computational Complexity}
\label{sec:method:complexity}

Let $N_t$ be the population size, $N_v=O(N_t)$ the number of directions, $N_c$
the size of the deduplicated candidate pool, $D$ the number of decision
variables and $M$ the number of objectives.

In the stratification module, direction association and the continuous PBI
computation of \eqref{eq:assoc-v}--\eqref{eq:score-cont} cost $O(N_tN_vM)$. By
\eqref{eq:direction-collapsed} the direction field requires only one normalisation
of the $n_{\mathrm{ND}}$ nondominated vectors, that is $O(n_{\mathrm{ND}}M)$; the
present implementation still calls $K$-means ($R=5$ replicates, iteration limit
$I=100$) at a cost of $O(R\,I\,n_{\mathrm{ND}}K_vM)$, which equals
$O(R\,I\,N_t^2M)$ once $K_v=n_{\mathrm{ND}}=O(N_t)$. This term is of the same order
as the dominating term of the module but carries a large constant factor; it is a
removable implementation overhead rather than an intrinsic cost of the method, and
Section~\ref{sec:exp:sensitivity} reports its share of the total runtime
separately. The
bounded bisection of the binary branch costs
$O(B_\delta N_t\keff M)$, where $B_\delta$ is fixed by the interval $[-20,20]$
and the tolerance $10^{-1}$, and anchor selection additionally performs
nondominated sorting and radar-projection distance computations. With
$N_v=O(N_t)$, $\keff=O(N_t)$ and fixed clustering limits, the module is dominated
by $O(N_t^2M)$; generating the relation pairs adds $O(N_t^2D)$.

In the candidate module, scoring one candidate requires $2N_t$ network inputs of
dimension $2D$, so the whole pool costs $O(N_cN_tD)$ in data construction and
storage, which dominates the module. The SDE fitness of \eqref{eq:sde} costs
$O(N_t^2M)$, the indicator mode adds one SVR fit on the $N_t$ evaluated solutions
plus predictions on the prefiltered candidates, and the greedy construction of
\eqref{eq:greedy} costs $O(n_{\max}N_cD)$. Overall the per-iteration cost is
therefore $O(N_cN_tD+N_t^2M)$.

Neither module consumes expensive objective evaluations: the total number of real
evaluations remains $B$ regardless of the mode sequence. The wall-clock share of
the two modules is reported in Section~\ref{sec:exp:sensitivity}.

Table~\ref{tab:params} lists the parameters introduced above.

\begin{table}[t]
\centering
\caption{Parameters of HPDC-MaOEA and their settings.}
\label{tab:params}
\small
\begin{tabular}{llc}
\toprule
Symbol & Meaning & Setting \\
\midrule
$\theta$ & penalty of the continuous PBI value & $5$ \\
$N_v$ & number of directions & $N_{\mathrm{init}}$ \\
$r_g$ & positive-group ratio & $0.25$ \\
$\delta$ & class-balance parameter of the binary branch &
  bisected in $[-20,20]$ \\
$\keff$ & effective anchor capacity & $\min(N_p,\max(6,\lceil1.5M\rceil))$ \\
$\alpha$ & weight of the continuous branch & $1-FE/B$ \\
$g_{\max}$ & candidate-generation limit & $3000$ \\
$q_{\mathrm{keep}}$ & quantile of the exploratory prefilter & $0.80$ \\
$\lambda_0$ & initial exploration strength & $0.35$ \\
$p_{\mathrm{mix}}$ & probability of the indicator-guided mode & $0.50$ \\
$n_{\min},n_{\max}$ & bounds of the evaluation batch & $4,6$ \\
\bottomrule
\end{tabular}
\end{table}

\section{Experimental Studies}
\label{sec:exp}

\subsection{Experimental Settings}
\label{sec:exp:settings}
\TODO{Benchmarks, $M\in\{3,5,10,20\}$, $D$, $N_p$, budget $B$, number of
independent runs, indicators (IGD/IGD$^+$/HV), statistical tests with Holm
correction, and platform. Each experiment carries exactly one argumentative
responsibility, per the plan in the problem-mechanism-evidence note (hybrid
PBI), Section 7.}

\subsection{Comparison with State-of-the-Art Algorithms}
\label{sec:exp:comparison}
\TODO{Algorithm-level performance against REMO, PIEA, PC-SAEA, R2AEA, RSEA and
further strong baselines under a common budget. This subsection establishes the
final performance difference only; the mechanism is isolated in
Sections~\ref{sec:exp:hpc} and~\ref{sec:exp:candidate}.}

\subsection{Ablation Study}
\label{sec:exp:ablation}
\TODO{Variants sharing the same host algorithm and the same $\keff$: continuous
branch only, binary branch only, and the full hybrid stratification; relation
top-$6$ selection versus dual-mode selection. Note the interaction between
$\keff$ and the label construction, and separate the anchor count from the radial
grid resolution of \eqref{eq:ndiv}.}

\subsection{Analysis of Hybrid PBI-Based Quality Stratification}
\label{sec:exp:hpc}
\TODO{Three mechanism questions, each with its own metric. (i) Non-redundancy of
the two views: mean Jaccard $0.2326$ over 250 runs; bidirectional unique future
true positives supported in 33 of 40 Problem--$M$--Stage cells after Holm
correction; $81.5\%$ of the union true positives unique to one view.
(ii) Source utilisation of the fused group: $32.2\%$ of the positive group from
the V-only region and $19.6\%$ from the A-only region, with $36.0\%$ and
$17.9\%$ of the future true positives. (iii) Budget-aware schedule versus static
and reversed schedules on a shared frozen candidate pool, plus the realised
$\alpha$ range of Remark~\ref{rem:alphafirst}.}

\subsection{Analysis of Dual-Mode Candidate Selection}
\label{sec:exp:candidate}
\TODO{Cutoff audit (400 runs: $5.1$-fold variation of the passing proportion,
$3.1$-fold variation of the iteration count) and the candidate-level probe over
four 20-objective problems with ten runs each: late-stage survival rate from
$0.706$ to $0.770$ (paired Wilcoxon, $p=0.0085$) and late-stage gain ratio
relative to a within-pool greedy reference from $0.0861$ to $0.1048$
($p=0.0250$). Report batch spread of each mode separately.}

\subsection{Parameter Sensitivity}
\label{sec:exp:sensitivity}
\TODO{Sensitivity to $r_g$, $q_{\mathrm{keep}}$, $\lambda_0$, $p_{\mathrm{mix}}$
and $[n_{\min},n_{\max}]$, and the wall-clock share of the two modules promised
in Section~\ref{sec:method:complexity}.}

\section{Conclusion}
\label{sec:conclusion}
\TODO{Conclusion. Reserved memory sentence: hybrid PBI stratification restores
the within-group ordering that binary partitioning compresses while retaining a
clear anchor-based boundary, and dual-mode candidate selection converts surrogate
scores into a bounded, scale-insensitive evaluation batch.}

\begin{thebibliography}{9}

\bibitem{ref:moead}
Q.~Zhang and H.~Li,
``MOEA/D: A multiobjective evolutionary algorithm based on decomposition,''
\emph{IEEE Trans. Evol. Comput.}, vol.~11, no.~6, pp.~712--731, 2007.

\bibitem{ref:remo}
\TODO{REMO -- relation-model-based expensive multi-objective optimization.}

\bibitem{ref:rsea}
\TODO{RSEA -- radar-grid based selection evolutionary algorithm.}

\bibitem{ref:piea}
Y.~Li, W.~Li, S.~Li, and Y.~Zhao,
``PIEA,'' \emph{Information Sciences}, 2024.
\TODO{complete volume and pages.}

\bibitem{ref:sde}
M.~Li, S.~Yang, and X.~Liu,
``Shift-based density estimation for Pareto-based algorithms in
many-objective optimization,''
\emph{IEEE Trans. Evol. Comput.}, vol.~18, no.~3, pp.~348--365, 2014.

\bibitem{ref:pcsaea}
\TODO{PC-SAEA -- pairwise-comparison-based surrogate-assisted EA,
Swarm Evol. Comput., 2023, doi:10.1016/j.swevo.2023.101323.}

\bibitem{ref:r2aea}
\TODO{R2AEA -- two-stage relation-based expensive optimization.}

\bibitem{ref:platemo}
Y.~Tian, R.~Cheng, X.~Zhang, and Y.~Jin,
``PlatEMO: A MATLAB platform for evolutionary multi-objective optimization,''
\emph{IEEE Comput. Intell. Mag.}, vol.~12, no.~4, pp.~73--87, 2017.

\end{thebibliography}

\end{document}
